\documentclass[11pt,a4paper]{article}

\usepackage[utf8]{inputenc}
\usepackage[T1]{fontenc}
\usepackage[margin=2.5cm]{geometry}
\usepackage{amsmath,amssymb}
\usepackage{booktabs}
\usepackage[hidelinks]{hyperref}

\title{Retrieval-Augmented Generation:\\
       Architecture, Failure Modes, and the Evaluation Problem}
\author{Prepared with LaTeXEditor}
\date{July 2026}

\begin{document}
\maketitle

\begin{abstract}
Retrieval-augmented generation conditions a language model on documents
fetched at inference time, decoupling knowledge from parameters. The
motivation is sound: knowledge stored in weights is expensive to update,
impossible to attribute, and prone to confident fabrication. This paper sets
out the architecture, then examines where deployed systems actually fail. We
argue that the field's attention has been disproportionately directed at the
retriever, while the dominant failure modes in practice arise at the boundary
between retrieval and generation --- the model ignoring retrieved evidence,
being actively degraded by irrelevant passages, or attributing a claim to a
document that does not support it. We also examine the chunking decision,
which is made early, rarely revisited, and constrains everything downstream.
Finally we address evaluation, where the compound nature of the system makes
end-to-end metrics uninformative about which component to fix.
\end{abstract}

\tableofcontents

\section{Introduction}
\label{sec:intro}

A language model trained on a corpus stores what it learned in its parameters.
This is a remarkable form of compression and a poor form of database. The
knowledge is fixed at training time, cannot be updated without further
training, cannot be attributed to a source, and --- although models do store a surprising amount of
relational fact \cite{petroni2019} --- cannot be distinguished by the
model from a plausible confabulation --- the mechanism that produces a correct
recollection is the same one that produces a fluent invention.

Retrieval-augmented generation \cite{lewis2020} addresses this by separating
the two functions. Knowledge lives in an external corpus; the model provides
language competence and reasoning. At inference the system retrieves relevant
passages and conditions generation on them. Updating knowledge becomes a
matter of updating documents. Attribution becomes possible because the
passages used are known. And the model is given evidence rather than being
asked to recall.

The approach is now the default for question answering over private corpora,
and a very large practical literature has grown around it \cite{gao2023}. Much of that
literature concerns retrieval quality. Our contention is that this emphasis is
misplaced: in deployed systems the retriever is usually adequate and the
failures occur elsewhere, in how the generator uses what it is given.

\section{Architecture}
\label{sec:architecture}

\subsection{The pipeline}

A minimal system has four stages. An offline stage segments documents into
passages and embeds them into a vector index. At query time the question is
embedded, the $k$ nearest passages retrieved, those passages placed in the
prompt, and an answer generated.

Formally, with query $q$, corpus $\mathcal{C}$ and retriever $R$,
\begin{equation}
  \hat{y} = \arg\max_y \ p_\theta\big(y \mid q, R_k(q, \mathcal{C})\big),
\end{equation}
where $R_k$ returns the $k$ highest-scoring passages. The original formulation
\cite{lewis2020} marginalises over retrieved documents,
\begin{equation}
  p(y \mid q) = \sum_{d \in R_k(q)} p_\eta(d \mid q) \, p_\theta(y \mid q, d),
  \label{eq:marginal}
\end{equation}
which permits training the retriever and generator jointly; REALM
\cite{guu2020} goes further and trains retrieval into pretraining itself. Almost all
deployed systems abandon this in favour of simple concatenation, because
Equation~\ref{eq:marginal} requires backpropagating into the retriever and
periodically re-indexing the corpus --- operationally awkward and rarely worth
it.

\subsection{Dense and sparse retrieval}

Sparse retrieval scores by lexical overlap. BM25 \cite{robertson2009} remains
the standard:
\begin{equation}
  \mathrm{BM25}(q, d) = \sum_{t \in q} \mathrm{IDF}(t)
    \cdot \frac{f(t, d)\,(k_1 + 1)}
                {f(t, d) + k_1\big(1 - b + b \frac{|d|}{\mathrm{avgdl}}\big)}.
\end{equation}
It requires no training, generalises across domains, and handles rare terms ---
product codes, names, identifiers --- that dense methods embed poorly.

Dense retrieval \cite{karpukhin2020} embeds queries and passages into a shared
space and scores by inner product, trained contrastively so that a question
lands near its answer passage. This captures paraphrase and synonymy that
lexical matching misses.

Neither dominates. Dense methods win on semantic similarity, sparse methods on
exact terms and out-of-domain corpora \cite{thakur2021}. Hybrid retrieval,
combining both scores, is standard in production for the straightforward
reason that their failures are uncorrelated.

\subsection{Approximate search}

Exhaustive inner-product search is linear in corpus size and infeasible at
scale, so indexes approximate. HNSW \cite{malkov2020} builds a navigable
small-world graph giving logarithmic search; inverted file indexes with
product quantisation \cite{johnson2019} partition the space and compress
vectors. Both trade recall for speed, and the recall lost is
\emph{invisible downstream}: the generator receives passages and cannot know
that a better one existed and was missed.

\begin{table}[htbp]
  \centering
  \caption{Retrieval approaches and their characteristic strengths.}
  \label{tab:retrieval}
  \begin{tabular}{llll}
    \toprule
    \textbf{Method} & \textbf{Matches} & \textbf{Strong on} & \textbf{Weak on} \\
    \midrule
    BM25 \cite{robertson2009}      & lexical  & rare terms, new domains & paraphrase \\
    Dense \cite{karpukhin2020}     & semantic & synonymy, intent        & exact identifiers \\
    Hybrid                          & both     & general use             & tuning the mixture \\
    Re-ranking \cite{nogueira2019} & pairwise & precision at small $k$  & latency \\
    \bottomrule
  \end{tabular}
\end{table}

\section{Chunking: the decision nobody revisits}
\label{sec:chunking}

Before any retrieval occurs, documents must be split. This choice is made
early, encoded in the index, and expensive to change --- and it constrains
everything downstream more than any subsequent component.

The tension is straightforward. Small chunks embed precisely, because a short
passage has a focused topic and its embedding is not an average of several.
But small chunks lack context: a paragraph beginning ``This approach fails
when the input is sparse'' is uninterpretable without knowing which approach.
Large chunks preserve context and embed diffusely, so their vectors sit near
the centroid of the space and retrieve poorly for specific questions.

Several mitigations are in use. Overlapping windows ensure no answer is split
across a boundary. Structural chunking respects document sections rather than
fixed token counts, which suits well-formatted sources and fails on
inconsistent ones. Contextual augmentation prepends a generated summary of the
parent document to each chunk before embedding, which addresses the
interpretability problem at the cost of an extra generation pass over the
corpus. Small-to-big retrieval embeds small chunks for precision but returns
their enclosing sections for context, decoupling the two requirements --- in
our judgement the most robust of the available options.

The general point deserves emphasis: a system whose retrieval is judged
inadequate frequently has a chunking problem rather than a retriever problem,
and swapping the embedding model will not fix it.

\section{Where deployed systems fail}
\label{sec:failures}

Assume retrieval works and the relevant passage is in the prompt. Several
failure modes remain, and these dominate in practice.

\subsection{The model ignores the evidence}

A model with strong parametric knowledge may answer from memory rather than
from the retrieved passage, particularly when they conflict. If the corpus has
been updated and the parameters have not, the model may confidently give the
outdated answer it was trained on while a correct passage sits in its context.

The behaviour depends on the strength of the prior. For obscure facts the
model defers to the context readily; for facts it is confident about, the
retrieved passage may be overridden. This is precisely the wrong way round for
a system whose purpose is to supply current information.

\subsection{Irrelevant context degrades performance}

More damaging, and less intuitive, is that retrieved passages which are
irrelevant do not merely fail to help --- they actively hurt
\cite{shi2023}. A model given a correct passage plus several irrelevant ones
performs worse than the same model given the correct passage alone.

This inverts the usual tuning instinct. Increasing $k$ raises the probability
that the answer is somewhere in the context, and simultaneously raises the
probability of distraction. Precision matters more than recall beyond a
modest $k$, which is the main argument for a re-ranking stage
\cite{nogueira2019}: retrieve broadly, then use a more expensive
cross-encoder to select the few passages actually worth including.

\subsection{Position effects}

Where a passage sits in the context matters. Models attend most reliably to
material at the beginning and end of their input and least reliably to the
middle \cite{liu2024} --- the ``lost in the middle'' effect. A system
retrieving ten passages and placing the best in the middle by relevance order
is placing it where it is least likely to be used. Ordering retrieved passages
so the highest-scoring sit at the extremes is a trivial change that measurably
improves accuracy, and is frequently omitted.

\subsection{Attribution that does not hold}

Systems that cite sources are trusted more, and the citation is often
generated rather than derived. A model asked to answer and cite may produce a
correct answer attributed to the wrong passage, or a claim that the cited
passage does not support. Because the citation looks like grounding, this is
more dangerous than an unattributed answer: it transfers unearned confidence.
Verifying that each generated claim is entailed by its cited passage requires a
separate entailment check, which few systems perform.

\section{Architectural variants}
\label{sec:variants}

Several designs depart from prompt concatenation.

\paragraph{Fusion-in-decoder.}
Rather than concatenating passages into one context, each is encoded
separately and the decoder attends across all of them \cite{izacard2021}. This
scales linearly in the number of passages rather than quadratically, allowing
many more, and avoids the position effects of Section~\ref{sec:failures}
because passages are not ordered in a sequence. It requires architectural
access to the model, so it is unavailable behind a generation API.

\paragraph{Retrieval at the token level.}
$k$NN-LM \cite{khandelwal2020} interpolates the model's next-token
distribution with one derived from nearest neighbours in a datastore of
training contexts. RETRO \cite{borgeaud2022} retrieves chunks during
pretraining and attends to them through dedicated cross-attention layers,
showing that a smaller model with retrieval can match a much larger one
without. These integrate retrieval more deeply than prompt-stuffing and
require training the model accordingly.

\paragraph{Adaptive retrieval.}
Not every query needs retrieval, and retrieving for a query the model can
answer directly adds latency and distraction risk. Self-RAG \cite{asai2024}
trains the model to emit control tokens deciding when to retrieve and to
critique whether retrieved passages are relevant and whether its own output is
supported. This makes the decision learned rather than fixed and provides a
natural place to hang a verification step.

\paragraph{Query transformation.}
User questions are frequently poor retrieval queries --- underspecified,
conversational, or containing pronouns referring to earlier turns. Rewriting
the query before retrieval, decomposing a multi-part question into
sub-questions, or generating a hypothetical answer and retrieving against that
all improve recall. These are cheap and effective, and in our experience yield
larger gains than changing the embedding model.

\section{Evaluation}
\label{sec:evaluation}

RAG is a compound system, and end-to-end accuracy conflates the contributions
of its parts. A wrong answer may reflect a chunking failure, a retrieval miss,
an approximate-index miss, distraction by irrelevant context, or a generation
error --- and the metric cannot distinguish them.

\subsection{Component metrics}

Retrieval is measured by recall@$k$ --- whether a passage containing the
answer appears in the top $k$ --- and by ranking measures such as MRR or
nDCG. These require relevance labels, which are expensive; BEIR
\cite{thakur2021} provides a standardised heterogeneous suite and is the usual
reference for retriever comparison.

Generation given context is measured by faithfulness (is every claim supported
by the provided passages?) and answer relevance (does it address the
question?). Faithfulness is the property that matters most and is the hardest
to measure automatically, since it requires entailment judgements over free
text.

\subsection{Model-based evaluation}

Frameworks such as RAGAS \cite{es2024} use a language model as judge,
decomposing an answer into claims and checking each against the context. This
scales where human annotation does not, and it inherits the judge's biases ---
including a documented preference for longer and more fluent answers
regardless of correctness. Model-based scores are useful for regression
detection and unreliable as absolute measures.

\subsection{What to measure in practice}

We would suggest three separable checks. First, retrieval recall against a
labelled set of query--passage pairs drawn from the actual corpus, which
isolates chunking and retrieval. Second, generation faithfulness given
\emph{gold} passages, which isolates the generator. Third, end-to-end accuracy,
which measures the system. Movement in the third without movement in the first
two indicates an interaction effect --- usually distraction --- and is
otherwise attributable.

\begin{table}[htbp]
  \centering
  \caption{Diagnosing a wrong answer by component.}
  \label{tab:diagnosis}
  \begin{tabular}{lll}
    \toprule
    \textbf{Symptom} & \textbf{Likely cause} & \textbf{Check} \\
    \midrule
    Answer absent from all passages & retrieval or chunking & recall@$k$ \\
    Correct passage present, ignored & distraction or position & re-rank, reorder \\
    Claim unsupported by cited passage & generation & faithfulness on gold \\
    Correct but outdated              & parametric override    & conflict test \\
    \bottomrule
  \end{tabular}
\end{table}

\section{Operational considerations}
\label{sec:operations}

\paragraph{Index freshness.}
The premise of RAG is that knowledge can be updated without retraining, which
holds only if the index is actually updated. Incremental indexing, deletion of
superseded documents, and detection of stale content are unglamorous and
determine whether the system delivers its central benefit.

\paragraph{Access control.}
In enterprise deployments not every user may see every document. Filtering
must occur at retrieval, not after generation --- a passage the user is not
entitled to see must never enter the prompt, because once it does, its content
can surface in the answer regardless of any post-hoc filter. Retrofitting
permissions onto an index built without them is difficult, which argues for
designing them in from the start.

\paragraph{Latency.}
Retrieval adds a round trip; re-ranking adds a model call; query rewriting adds
another. A pipeline with all three can spend more time before generation than
generating. Where latency matters, adaptive retrieval that skips the pipeline
for queries not requiring it is the most effective lever.

\paragraph{Cost.}
Retrieved context dominates prompt length, and prompt length dominates cost.
Reducing $k$ from ten passages to three often improves accuracy, by the
distraction argument, while cutting cost proportionally --- one of the few
genuinely free improvements available.

\section{Conversational retrieval}
\label{sec:conversational}

Most deployed systems are not single-turn. A user asks a follow-up, and the
follow-up is frequently unintelligible in isolation --- ``what about the
second one?'', ``does that apply to contractors?''. Retrieval against such a
query returns nothing useful, because the terms that would make it specific
are in the preceding turns.

The standard remedy is query rewriting: a model reformulates the follow-up
into a self-contained question using the conversation history, and retrieval
runs against the rewrite. This works well and introduces a failure mode of its
own, since the rewrite may resolve a reference incorrectly and send retrieval
confidently in the wrong direction. A rewrite is an unlogged intermediate in
most implementations, which makes such failures hard to diagnose --- the
system appears to retrieve badly for a question the user never asked.

An alternative is to retrieve against the concatenated history, which avoids
the intermediate but dilutes the query with earlier topics. A third is to
maintain a running summary of the conversation state and retrieve against
both it and the current turn. We would recommend logging the rewritten query
regardless of approach; it is the cheapest available observability measure and
resolves a large fraction of otherwise mysterious retrieval failures.

\paragraph{History and distraction.}
The distraction result of Section~\ref{sec:failures} applies to conversation
history as well as retrieved passages. A long history of prior turns occupies
context and competes for attention with the retrieved evidence. Systems that
accumulate turns indefinitely degrade for this reason, and truncating or
summarising history is not merely a cost measure.

\section{Structured and hybrid sources}
\label{sec:structured}

The architecture of Section~\ref{sec:architecture} assumes the corpus is
unstructured text. Many real corpora are not, and forcing them into that shape
loses information.

\paragraph{Tables.}
A table serialised into a text chunk embeds poorly: the vector reflects the
column headers and general topic, not the specific values. A question asking
for a particular cell is unlikely to retrieve the table containing it. Better
results come from generating a natural-language description of each table for
retrieval while returning the table itself for generation --- the
small-to-big pattern of Section~\ref{sec:chunking} applied to a different
mismatch.

\paragraph{Relational data.}
Where the answer requires aggregation --- how many, what is the average,
which is largest --- retrieval cannot supply it, because no passage contains
the answer. The correct architecture generates a query against a database
rather than retrieving text. Systems that attempt aggregation through
retrieval fail in a characteristic way: they retrieve several relevant records
and produce a confident count based on those, which is wrong whenever the true
count exceeds $k$.

\paragraph{Knowledge graphs.}
Where entities and their relationships are explicit, graph traversal answers
multi-hop questions that flat retrieval handles badly. Retrieving for ``who
succeeded the person who founded this company'' requires two hops, and a
single dense query is unlikely to surface both. Hybrid systems retrieve an
entity subgraph and pass its serialisation to the generator. The cost is that
the graph must exist, and constructing one from unstructured text is itself
error-prone.

The general lesson is that ``retrieval'' is the right primitive for questions
answered by a passage, and the wrong one for questions answered by a
computation over many records. Recognising which kind a query is --- and
routing accordingly --- is a more valuable capability than improving either
path in isolation.

\section{Security considerations}
\label{sec:security}

Conditioning generation on retrieved documents means allowing documents to
influence behaviour, and where those documents are not fully trusted this is a
vulnerability.

\paragraph{Indirect prompt injection.}
A document in the corpus may contain text addressed to the model rather than
the reader --- instructions to ignore prior guidance, to exfiltrate context,
or to produce a particular answer. Because retrieved content enters the same
context as the system prompt, the model has no reliable way to distinguish
data from instruction. In a corpus that any user can contribute to, or that
ingests external web content, this is a live risk rather than a theoretical
one.

Mitigations are partial. Delimiting retrieved content and instructing the
model to treat it as data raises the bar without settling the matter.
Sanitising ingested documents for instruction-like patterns catches obvious
cases. The structural defence is to constrain what the system can do with a
compromised generation --- no tool calls derived from retrieved instructions,
no outbound requests to URLs found in passages --- so that injection cannot
convert into action.

\paragraph{Corpus poisoning.}
An adversary able to insert documents can craft passages engineered to be
retrieved for particular queries and to contain misleading content. Dense
retrievers are susceptible to embedding-space attacks, where a passage is
optimised to sit near a target query vector while reading naturally. Provenance
tracking and restricting write access to the index are the practical defences.

\paragraph{Extraction of the corpus.}
A system that faithfully quotes retrieved passages will, given enough
carefully chosen queries, disclose substantial portions of a private corpus.
Where the corpus is confidential this is a data-loss channel that the access
controls of Section~\ref{sec:operations} address only if enforced at retrieval
time.

\section{When not to use retrieval}
\label{sec:whennot}

A survey of an approach should say where it does not apply, and RAG is applied
in several places it does not suit.

\paragraph{Stable, compact knowledge.}
If the knowledge is small, unchanging and central to every query, putting it in
the system prompt is simpler, faster and more reliable than retrieving it.
Retrieval introduces a failure mode --- the passage is not retrieved --- that
a fixed prompt does not have.

\paragraph{Reasoning-dominated tasks.}
Where the difficulty is inference rather than recall, retrieval adds context
without adding capability, and by the distraction argument may reduce
performance. Mathematical problem solving is the clearest case: the relevant
knowledge is already parametric and the failure is in the reasoning chain.

\paragraph{Style and format tasks.}
Asking a model to rewrite text in a house style is better served by
fine-tuning or few-shot examples than by retrieving style guidelines, because
the requirement is procedural rather than factual. Retrieved rules describe
what to do; examples demonstrate it, and demonstration transfers better.

\paragraph{Genuinely novel synthesis.}
Where no document contains the answer because the answer does not yet exist,
retrieval supplies inputs at best. Presenting such a system as retrieval-backed
overstates the grounding: the synthesis is the model's, and the citations
support the premises rather than the conclusion.

The practical recommendation is to route rather than to apply RAG uniformly.
Classifying an incoming query by what it needs --- fixed context, retrieval,
computation, or none --- and dispatching accordingly produces a better system
than any single path, and the classifier need not be sophisticated to beat
always-retrieve.

\section{Discussion}
\label{sec:discussion}

The framing that has served us best is that RAG is not a technique but a
system, and its behaviour is determined by interactions between components
rather than by the quality of any one.

The evidence for this is that the highest-leverage interventions we have
described are not retriever improvements. Reordering passages so the best sit
at the context extremes; reducing $k$ to limit distraction; chunking
small-to-big; rewriting queries before retrieval; checking entailment before
citing. None involves a better embedding model, and each addresses a boundary
between stages.

This suggests that the research emphasis on retrieval quality --- where
benchmarks are clean and improvements measurable --- has been partly a
streetlight effect. The failures that matter are harder to benchmark because
they are properties of the composition.

\section{Conclusion}
\label{sec:conclusion}

Retrieval-augmented generation separates knowledge from language competence,
making knowledge updatable and answers attributable. Both benefits are real
and neither is automatic.

Deployed systems fail predominantly at the retrieval--generation boundary
rather than within retrieval: models override provided evidence with
parametric knowledge, degrade in the presence of irrelevant passages, attend
unevenly across context positions, and generate citations that do not support
their claims. The chunking decision, made before any of this and rarely
revisited, constrains what retrieval can achieve at all.

Evaluation should therefore be componentwise. End-to-end accuracy tells a team
that something is wrong; separate measurement of retrieval recall and
grounded-generation faithfulness tells them what. Systems reporting only the
former are difficult to improve deliberately, and improvements that do arrive
are difficult to attribute.

\begin{thebibliography}{99}

\bibitem{asai2024}
A.~Asai, Z.~Wu, Y.~Wang, A.~Sil, and H.~Hajishirzi.
\newblock Self-RAG: Learning to retrieve, generate, and critique through
  self-reflection.
\newblock In \emph{ICLR}, 2024.

\bibitem{borgeaud2022}
S.~Borgeaud et~al.
\newblock Improving language models by retrieving from trillions of tokens.
\newblock In \emph{ICML}, 2022.

\bibitem{es2024}
S.~Es, J.~James, L.~Espinosa-Anke, and S.~Schockaert.
\newblock RAGAS: Automated evaluation of retrieval augmented generation.
\newblock In \emph{EACL}, 2024.

\bibitem{gao2023}
Y.~Gao, Y.~Xiong, X.~Gao, K.~Jia, J.~Pan, Y.~Bi, Y.~Dai, J.~Sun, and H.~Wang.
\newblock Retrieval-augmented generation for large language models: A survey.
\newblock \emph{arXiv:2312.10997}, 2023.

\bibitem{guu2020}
K.~Guu, K.~Lee, Z.~Tung, P.~Pasupat, and M.-W. Chang.
\newblock REALM: Retrieval-augmented language model pre-training.
\newblock In \emph{ICML}, 2020.

\bibitem{izacard2021}
G.~Izacard and E.~Grave.
\newblock Leveraging passage retrieval with generative models for open domain
  question answering.
\newblock In \emph{EACL}, 2021.

\bibitem{johnson2019}
J.~Johnson, M.~Douze, and H.~Jégou.
\newblock Billion-scale similarity search with GPUs.
\newblock \emph{IEEE Transactions on Big Data}, 7(3):535--547, 2019.

\bibitem{karpukhin2020}
V.~Karpukhin, B.~Oğuz, S.~Min, P.~Lewis, L.~Wu, S.~Edunov, D.~Chen, and
  W.-t. Yih.
\newblock Dense passage retrieval for open-domain question answering.
\newblock In \emph{EMNLP}, 2020.

\bibitem{khandelwal2020}
U.~Khandelwal, O.~Levy, D.~Jurafsky, L.~Zettlemoyer, and M.~Lewis.
\newblock Generalization through memorization: Nearest neighbor language
  models.
\newblock In \emph{ICLR}, 2020.

\bibitem{lewis2020}
P.~Lewis et~al.
\newblock Retrieval-augmented generation for knowledge-intensive NLP tasks.
\newblock In \emph{NeurIPS}, 2020.

\bibitem{liu2024}
N.~F. Liu, K.~Lin, J.~Hewitt, A.~Paranjape, M.~Bevilacqua, F.~Petroni, and
  P.~Liang.
\newblock Lost in the middle: How language models use long contexts.
\newblock \emph{TACL}, 12:157--173, 2024.

\bibitem{malkov2020}
Y.~A. Malkov and D.~A. Yashunin.
\newblock Efficient and robust approximate nearest neighbor search using
  hierarchical navigable small world graphs.
\newblock \emph{IEEE TPAMI}, 42(4):824--836, 2020.

\bibitem{nogueira2019}
R.~Nogueira and K.~Cho.
\newblock Passage re-ranking with BERT.
\newblock \emph{arXiv:1901.04085}, 2019.

\bibitem{petroni2019}
F.~Petroni, T.~Rocktäschel, P.~Lewis, A.~Bakhtin, Y.~Wu, A.~H. Miller, and
  S.~Riedel.
\newblock Language models as knowledge bases?
\newblock In \emph{EMNLP}, 2019.

\bibitem{robertson2009}
S.~Robertson and H.~Zaragoza.
\newblock The probabilistic relevance framework: BM25 and beyond.
\newblock \emph{Foundations and Trends in Information Retrieval},
  3(4):333--389, 2009.

\bibitem{shi2023}
F.~Shi, X.~Chen, K.~Misra, N.~Scales, D.~Dohan, E.~Chi, N.~Schärli, and
  D.~Zhou.
\newblock Large language models can be easily distracted by irrelevant
  context.
\newblock In \emph{ICML}, 2023.

\bibitem{thakur2021}
N.~Thakur, N.~Reimers, A.~Rücklé, A.~Srivastava, and I.~Gurevych.
\newblock BEIR: A heterogeneous benchmark for zero-shot evaluation of
  information retrieval models.
\newblock In \emph{NeurIPS Datasets and Benchmarks}, 2021.

\end{thebibliography}

\end{document}
